
\section{Conclusion}

\subsection{Summary}

We presented \textbf{RPC} (Random Projection Compressed) attention, which sketches 
$Q$ and $K$ into a lower-dimensional subspace before score computation. 
By the Johnson--Lindenstrauss lemma, the sketched scores approximate the 
full scores, reducing attention FLOPs.

\subsection{Results Recap}

\begin{table}[H]
\centering
\small
\begin{tabular}{llll}
\toprule
\textbf{Model} & \textbf{Acc $\Delta$} & \textbf{Param $\Delta$ \%} & \textbf{Time ratio (L=128)} \\
\midrule
DistilBERT & {{D\_DELTA}}\% & +{{D\_PARAM\_PCT}}\% & {{D\_TIME\_RATIO}}$\times$ \\
RoBERTa & {{R\_DELTA}}\% & +{{R\_PARAM\_PCT}}\% & {{R\_TIME\_RATIO}}$\times$ \\
GPT-2 & {{G\_DELTA}}\% & +{{G\_PARAM\_PCT}}\% & {{G\_TIME\_RATIO}}$\times$ \\
\bottomrule
\end{tabular}
\end{table}

\subsection{Key Takeaways}

\begin{enumerate}
\item \textbf{Not a compressor.} Parameters increase.
\item \textbf{Not a memory saver.} Score matrix unchanged.
\item \textbf{FLOPs reduction only at long sequences.}
\item \textbf{Theoretical speedup $\ne$ actual speedup.}
\item \textbf{Multi-architecture support works.}

\end{enumerate}
\subsection{Future Work}

\subsubsection{Short-Term}
\begin{itemize}
\item Test at L=1024, L=4096
\item Combine with AASVD
\item Sparse projections

\end{itemize}
\subsubsection{Long-Term}
\begin{itemize}
\item Learned projections with JL constraints
\item Theoretical bounds on downstream accuracy
\item Comparison with Performer

\end{itemize}
\subsection{Closing Remarks}

RPC provides a mathematically principled approach to reducing attention 
FLOPs. It is not a compressor and does not reduce memory. We release all 
code and checkpoints.

\subsection{Acknowledgments}

{{ACKNOWLEDGMENTS}}

\subsection{Author Contributions}

\begin{itemize}
\item \textbf{{{AUTHOR\_1}}} designed RPC, implemented the code, ran experiments,
  wrote the paper.
\item \textbf{{{AUTHOR\_2}}} provided feedback, edited the manuscript.

\end{itemize}
\subsection{Conflict of Interest}

None.

